\scanpage{136}

\unsure{éléments|élément} dont l'image dans $\mathcal E_K^{\bar K}$ est dans l'image de $\mathcal E_{k(x)}^{\overline{k(x)}}\unsure{\mathrel{\text{---}}|\longrightarrow}\mathcal E_K^{\bar K}$ --- on voit que pour $x$ variable dans une composante connexe $D_{J,\alpha}$ de $D_J$, sont tous conjugués entre eux dans $\mathcal E_{U/K}^{\unsure{\bar L|\bar K|\deleted{\bar K}}}$ --- et on peut faire les choix relatifs aux différents $\alpha$ d'une façon cohérente telle\unsure{.}, que maintenant les autom extérieurs de $\pi_{U/K}^{\bar L}$ qui expriment l'indétermination dans (39) soient ceux provenant des éléments de l'image de
\[
\mathcal E_{D_{J,\alpha}}^{\tilde D_{J,\alpha}}\longrightarrow\mathcal E_K^{\bar K}.
\]

$J=\{i\}$ est particulièrement \unsure{simple}, \deleted{chaque composante \unsure{conn.}}\inserted{dans ce cas d'ailleurs $D_J=D_i$} est \unsure{connëxe|connexe}, et on trouve des hom
\[
\tag{40}
T\xrightarrow{\kappa_i}\pi_{U/K}^{\bar L},
\]
déterminés a priori mod autom intérieurs près --- donc en fait $\kappa_i$ est un \uline{hom} \uline{extérieur} bien déterminé, qui est \deleted{\unsure{invariant}} \deleted{stable par l'opération} compatible avec les opérations ext. de $\mathcal E_K^{\bar K}$ sur $T$, et sur $\pi_{U/K}^{\bar L}$.

\margin{de façon plus précise, le choix d'un rev. universel de $\tilde U_x$ en implique un de $K$ --- et on suppose que c'est $\bar K$ (avec \unsure{des} \unsure{rigidifications|rigidification}) i.e. on choisit $\tilde U_x$ au dessus de $\bar K$ --- et alors cela implique un rev. universel de $k(x)$ au dessus de $\bar K$ \unsure{aussi}. Ces choix étant faits, on a
\diagram{
\begin{tikzcd}[ampersand replacement=\&]
\mathcal E_{\tilde U_x}^{\unsure{\tilde{\tilde U}_x}} \arrow[r] \& \mathcal E_K^{\bar K}
\end{tikzcd}
}
\ldots}

\margin{Mais on a mieux ils sont \unsure{(}par voie géom.) définis mod autom intérieurs de $\uline{\pi_{U/K}^{\bar L}}$, et les opérations \deleted{\unsure{extérieures}} de $\mathcal E_K^{\bar K}$ sur $\pi_{U/K}^{\bar L}$ laissent invariants les hom \inserted{extérieurs} en tenant compte de l'action extérieure ``cyclotomique'' de $\mathcal E_K^{\bar K}$ sur $T^J$\ldots}

\deleted{Si \unsure{on a} $J'\supset J$ \inserted{(\unsure{donc|d'où} $\bar D_{J'}\subset\bar D_J$),} (i.e. $\bar D_{J'}\deleted{\unsure{\cap D_J}}\ne\phi$), choisissons une composante conn. $\unsure{\bar D_{J',\alpha'}}$ de $D_{J'}$, elle est contenue dans une comp. conn.}

\margin{135}
